\documentclass[11pt,a4paper]{article}

\usepackage[margin=2.5cm]{geometry}
\usepackage{amsmath}
\usepackage{amssymb}
\usepackage{booktabs}
\usepackage{hyperref}
\usepackage{parskip}
\usepackage{microtype}
\usepackage{array}
\usepackage{cite}

\hypersetup{
    colorlinks=true,
    linkcolor=black,
    citecolor=black,
    urlcolor=blue
}

\title{\textbf{Nature-Inspired Optimisation Algorithms}\\
       \large A Literature Review: CITS4404 Part 1}
\author{James Wigfield (23334375) \and
        Farahnaz Abdi (23475614) \and
        Gaganish Yadav (24552039) \and
        Chenxu You (24547664)}
\date{April 2026}

\begin{document}

\maketitle

\section{Dragonfly Algorithm (DA) Synopsis}

\noindent\textbf{Source paper:} S.\ Mirjalili, ``Dragonfly algorithm: a new
meta-heuristic optimization technique for solving single-objective, discrete,
and multi-objective problems,'' \textit{Neural Computing and Applications},
vol.~27, pp.~1053--1073, 2016.~\cite{mirjalili2016dragonfly} 

\subsection*{Motivation and Gap}

The Dragonfly Algorithm (DA) is a swarm intelligence (SI) optimiser proposed by
Mirjalili~\cite{mirjalili2016dragonfly} to handle continuous, binary, and
multi-objective optimisation problems.
While algorithms like Particle Swarm Optimisation
(PSO)~\cite{kennedy1995pso}, Ant Colony Optimisation (ACO), and Artificial Bee
Colony (ABC) have been widely studied, they only model a small some of the 
behaviours seen in natural swarms, and dragonfly swarming in particular
had never been explored computationally before this paper.
The No Free Lunch theorem~\cite{wolpert1997nfl} also supports the case for a new
approach: since no single algorithm works best on every problem, a genuinely
different method still has room to outperform existing ones on the problems where
they struggle.

The specific gap DA targets is that previous SI algorithms model some aspects of
swarm behaviour but none capture all five survival behaviours at once.
For example, PSO only models attraction toward a personal best and a global
best, and while several improved variants have tried to add separation,
alignment, or cohesion operators
independently~\cite{cui2009boid,kadrovach2002particle,cui2009alignment}, none
brought all five together in a single framework.
Beyond that, most existing algorithms don't separate two quite different swarming
modes: \textit{static} swarms that hunt over a small area (exploration) and
\textit{dynamic} migratory swarms that travel in one direction (exploitation).
A further gap is the absence of an \textit{enemy} concept, search agents are
only pulled toward promising regions but never actively pushed away from bad
ones, which can hurt diversity early in the search.

\subsection*{The Algorithm}

The main idea is to model artificial dragonflies using \textbf{five behavioural
operators} built on Reynolds' flocking rules~\cite{reynolds1987flocks}, with two
extra terms added for food and enemy interactions:

\begin{enumerate}
    \item \textbf{Separation} ($S_i = -\sum_j (X - X_j)$): avoids collision with neighbours.
    \item \textbf{Alignment} ($A_i = \frac{1}{N}\sum_j V_j$): matches neighbours' velocities.
    \item \textbf{Cohesion} ($C_i = \frac{1}{N}\sum_j X_j - X$): moves toward the neighbourhood centre of mass.
    \item \textbf{Food attraction} ($F_i = X^+ - X$): steers toward the current best solution.
    \item \textbf{Enemy distraction} ($E_i = X^- + X$): steers away from the current worst solution.
\end{enumerate}

The \textbf{step vector} (analogous to PSO velocity) combines all five:
\[
    \Delta X_{t+1} = (s\,S_i + a\,A_i + c\,C_i + f\,F_i + e\,E_i) + w\,\Delta X_t
\]
where $w$ decays over iterations to shift from exploration to exploitation.
When a dragonfly has no neighbours it performs a \textbf{L\'{e}vy flight}
($X_{t+1} = X_t + \text{Levy}(d)\cdot X_t$), giving agents with no nearby neighbours a way to
still explore the space.
Static swarming (high alignment weight $a$, low cohesion weight $c$) models
exploration, dynamic swarming (low $a$, high $c$) models exploitation, and the
neighbourhood radius increases over time so what starts as many small independent
swarms gradually merges into a larger swarm.

There are two main variants that extend this base algorithm.
The \textbf{Binary DA (BDA)} uses a v-shaped transfer function to convert the
continuous step values into bit-flip probabilities, v-shaped functions work
better here because s-shaped ones tend to push particles toward extreme binary
values~\cite{saremi2014transfer}.
The \textbf{Multi-Objective DA (MODA)} adds a Pareto archive and a grid system
where food sources are picked from the least-crowded cells (to spread solutions
out) and enemies from the most-crowded ones (to avoid revisiting dense areas),
borrowing the archive approach from MOPSO~\cite{coello2004mopso}.

Each variant is tested separately against established baselines.
DA is benchmarked on 19 standard test functions (unimodal, multi-modal, and
composite) against PSO and GA, with the Wilcoxon rank-sum test used to check
whether differences are statistically significant.
BDA is tested on 13 of the same functions encoded in binary space, compared
against Binary PSO and Binary GSA.
MODA is evaluated on five ZDT multi-objective benchmarks plus a real submarine
propeller design problem, using Inverse Generalised Distance (IGD) as the
performance metric and comparing against MOPSO and NSGA-II.

% The paper includes pseudocode, equations, and parameter settings in enough
% detail to replicate the experiments, and source code is publicly available.

\subsection*{Results and Assessment}

On single-objective problems, DA achieves the best or joint-best result on 13 of
19 test functions.
The advantage over GA is clear and statistically significant on almost all
functions ($p < 0.0002$).
Against PSO the picture is more mixed, differences are often not significant
on multi-modal functions (e.g.\ TF9, TF11, TF12), and on TF8 (Schwefel) PSO
substantially outperforms DA.
On composite functions the advantage over PSO is not consistent, though the paper
puts this down to those functions just being harder to solve in general.
BDA outperforms BPSO and BGSA on 11 of 13 binary functions ($p < 0.002$), and
MODA beats NSGA-II on all five ZDT variants by a clear margin while being
competitive with MOPSO and outperforming it on ZDT3 and the three-objective variant.

Overall, the main claims hold up reasonably well.
The methodology is solid (30-run averaging, Wilcoxon significance testing, and
a broad benchmark set).
One thing worth noting is that DA is only compared against PSO and GA for the
single-objective case, there's no comparison against anything more recent like
Differential Evolution or the Grey Wolf Optimiser, so it's hard to know how well
it would stack up against the current state of the art.
The mixed results on composite functions also suggest the exploration/exploitation
balance doesn't always hold up on more complex landscapes.
For this assignment though, DA looks like a practical fit, its continuous
parameter space and the way the five operators naturally balance exploration and
exploitation make it well-suited to optimising the 7-dimensional trading bot
parameter vectors, and the pseudocode in the paper is clear enough to implement
from scratch.

% Algorithm: ..
\section{[Algorithm 2 Name] Synopsis}
\label{sec:algo2}

\noindent\textbf{Source paper:} 

\subsection*{Motivation and Gap}


\subsection*{The Algorithm}


\subsection*{Results and Assessment}

% Algorithm: 
\section{[Algorithm 3 Name] Synopsis}
\label{sec:algo3}

\noindent\textbf{Source paper:} 
\subsection*{Motivation and Gap}


\subsection*{The Algorithm}


\subsection*{Results and Assessment}


% Algorithm .
\section{[Algorithm 4 Name] Synopsis}
\label{sec:algo4}

\noindent\textbf{Source paper:} 
\subsection*{Motivation and Gap}

\subsection*{The Algorithm}


\subsection*{Results and Assessment}


\section{Comparative Conclusion}
\label{sec:conclusion}



\bibliographystyle{IEEEtran}
\begin{thebibliography}{99}

% DA references (James Wigfield)

\bibitem{mirjalili2016dragonfly}
S.~Mirjalili, ``Dragonfly algorithm: a new meta-heuristic optimization technique
for solving single-objective, discrete, and multi-objective problems,''
\textit{Neural Comput.\ Appl.}, vol.~27, no.~4, pp.~1053--1073, 2016.

\bibitem{kennedy1995pso}
J.~Kennedy and R.~Eberhart, ``Particle swarm optimization,'' in
\textit{Proc.\ IEEE Int.\ Conf.\ Neural Networks}, 1995, pp.~1942--1948.

\bibitem{wolpert1997nfl}
D.~H.~Wolpert and W.~G.~Macready, ``No free lunch theorems for optimization,''
\textit{IEEE Trans.\ Evol.\ Comput.}, vol.~1, no.~1, pp.~67--82, 1997.

\bibitem{reynolds1987flocks}
C.~W.~Reynolds, ``Flocks, herds and schools: a distributed behavioral model,''
\textit{ACM SIGGRAPH Comput.\ Graph.}, vol.~21, no.~4, pp.~25--34, 1987.

\bibitem{saremi2014transfer}
S.~Saremi, S.~Mirjalili, and A.~Lewis, ``How important is a transfer function
in discrete heuristic algorithms,''
\textit{Neural Comput.\ Appl.}, pp.~1--16, 2014.

\bibitem{coello2004mopso}
C.~A.~C.~Coello, G.~T.~Pulido, and M.~S.~Lechuga, ``Handling multiple
objectives with particle swarm optimization,''
\textit{IEEE Trans.\ Evol.\ Comput.}, vol.~8, no.~3, pp.~256--279, 2004.

\bibitem{cui2009boid}
Z.~Cui and Z.~Shi, ``Boid particle swarm optimisation,''
\textit{Int.\ J.\ Innov.\ Comput.\ Appl.}, vol.~2, pp.~77--85, 2009.

\bibitem{kadrovach2002particle}
B.~A.~Kadrovach and G.~B.~Lamont, ``A particle swarm model for swarm-based
networked sensor systems,'' in
\textit{Proc.\ ACM Symp.\ Appl.\ Comput.}, 2002, pp.~918--924.

\bibitem{cui2009alignment}
Z.~Cui, ``Alignment particle swarm optimization,'' in
\textit{Proc.\ 8th IEEE Int.\ Conf.\ Cognitive Informatics}, 2009,
pp.~497--501.

% Additional references:

\end{thebibliography}

\end{document}
